\documentclass[12pt, a4paper]{article}
\usepackage[utf8]{inputenc}
\usepackage[german]{babel}
\usepackage[T1]{fontenc}
\usepackage{amsmath, amssymb, amsthm}
\usepackage{graphicx}
\usepackage{geometry}
\usepackage{xcolor}
\usepackage{tikz}
\usepackage{pgfplots}

\geometry{top=2cm, bottom=2cm, left=2.5cm, right=2.5cm}
\pgfplotsset{compat=1.18}

% Titel und Header
\title{\textbf{Advanced KantiKoala LaTeX Test}}
\author{Demo Teacher}
\date{\today}

\begin{document}

\maketitle

\section{Einführung}
Dieses Dokument demonstriert, dass nun \textbf{alle} auf dem Server installierten LaTeX-Pakete unterstützt werden. Da das Dokument als echtes PDF kompiliert und eingebettet wird, ist das Layout 100\% originalgetreu.

\section{Paket-Unterstützung}

Die folgenden Pakete sind Beispiele, die in der einfachen HTML-Konvertierung nicht funktioniert hätten, jetzt aber möglich sind:

\subsection{TikZ Grafiken}
Hier ist eine Vektorgrafik, die direkt in LaTeX generiert wird:

\begin{center}
\begin{tikzpicture}
    \draw[fill=blue!20] (0,0) circle (1.5cm);
    \draw[->, thick, red] (-2,0) -- (2,0) node[right] {$x$};
    \draw[->, thick, red] (0,-2) -- (0,2) node[above] {$y$};
    \node at (0.5,0.5) {$r=1.5$};
    \draw (0,0) -- (1.06, 1.06);
\end{tikzpicture}
\end{center}

\subsection{PGFPlots Funktionsgraphen}
\begin{center}
\begin{tikzpicture}
\begin{axis}[
    axis lines = middle,
    xlabel = $x$,
    ylabel = {$f(x)$},
    width=8cm, height=6cm,
]
\addplot [
    domain=-2:2, 
    samples=100, 
    color=blue,
    thick,
]
{x^3 - x};
\addlegendentry{$x^3 - x$}
\end{axis}
\end{tikzpicture}
\end{center}

\section{Mathematik (Advanced)}
Natürlich funktioniert auch komplexe Mathematik weiterhin tadellos, da es natives LaTeX ist.

$$ \oint_C \vec{F} \cdot d\vec{r} = \iint_S (\nabla \times \vec{F}) \cdot d\vec{S} $$

\section{Hinweis für Lehrpersonen}
Welche Pakete können verwendet werden?
\begin{itemize}
    \item \textbf{Alle Standardpakete} (geometry, fancyhdr, babel, etc.)
    \item \textbf{TikZ / PGFPlots} für Diagramme
    \item \textbf{Tabularx / Longtable} für komplexe Tabellen
    \item \textbf{Hinweis:} Externe Bilder (JPG/PNG) funktionieren nur, wenn sie auf dem Server liegen (was schwierig ist). Es wird empfohlen, Grafiken mit TikZ zu erstellen oder Word zu nutzen, wenn viele Bilder benötigt werden.
\end{itemize}

\end{document}
